      下面给出陈述。
      断言这三个向量不共面，就等同于断言没有向量落在另外两个向量的张成之中\Dash 没有向量是另两个向量的线性组合。
      这恰恰就是断言这个三元集合是线性无关的。
       由\nearbycorollary{cor:NVecsRNSpanIffLI}可知，这等价于断言该集合是 $\Re^3$ 的一个基（更精确地说，由该集合的元素构成的任何序列都是一个基）。
    \end{answer}
  \item 
    % \begin{exparts}
      % \partsitem Prove that any subspace of a finite dimensional space
      %   has a basis. 
        证明有限维空间的任何子空间都是有限维的。
    % \end{exparts}
    \begin{answer}
      设空间 $V$ 是有限维的，且 $S$ 是 $V$ 的一个子空间。

      若 $S$ 不是有限维的，那么它有一个无穷的线性无关集（从空集出发，不断添加不与该集合线性相关的向量；这一过程可以持续无穷多步，否则 $S$ 就会是有限维的）。
      但 $S$ 的任何线性无关子集也是 $V$ 的线性无关子集，这与\nearbycorollary{cor:NoLiSetGreatDim}矛盾
      % \begin{exparts}
      %   \partsitem The empty set is a linearly independent subset of $S$.
      %     By \nearbycorollary{cor:LIExpBas}, it can be expanded to a basis
      %     for the vector space $S$. 
      %   \partsitem Any basis for the subspace $S$ is a linearly independent 
      %     set in the superspace $V$.
      %     Hence it can be expanded to a basis for the superspace, which is
      %     finite dimensional.
      %     Therefore it has only finitely many members.  
      % \end{exparts}
    \end{answer}
  \item  
    在\nearbytheorem{th:AllBasesSameSize}中，\( B \) 的有限性在哪里被用到？
    \begin{answer}
      它保证我们能穷尽 \( \vec{\beta} \) 中的各个向量。也就是说，它为最后一段的第一句话提供了依据。
    \end{answer}
  \item
    证明：若 \( U \) 与 \( W \) 都是 \( \Re^5 \) 的三维子空间，则 \( U\intersection W \) 是非平凡的。试加以推广。
    \begin{answer}
       设 \( B_U \) 是 \( U \) 的一个基，\( B_W \) 是 \( W \) 的一个基。
       考虑将这两个基的序列首尾相连。
       若有重复出现的元素，则交 \( U\intersection W \) 是非平凡的。
       否则，
       集合 \( B_U\union B_W \) 是线性相关的，因为它是五维空间 \( \Re^5 \) 的一个六元子集。
       无论哪种情形，\( B_W \) 中总有某个成员落在 \( B_U \) 的张成之中，于是 \( U\intersection W \) 就不只是平凡空间 \( \set{\zero\,} \)。

       推广：
       若 \( U,W \) 是维数为 \( n \) 的向量空间的子空间，且 \( \dim(U)+\dim(W)>n \)，则它们的交是非平凡的。
    \end{answer}
  \item 
    一个空间的基由该空间的元素组成。
    于是我们自然会问：“是基”这一性质与运算 $\subseteq$、$\cap$、$\cup$ 是如何相互作用的。
    （当然，基实际上是一个有序的序列，但这些运算有一种自然的推广。）
    \begin{exparts}
      \partsitem 先考虑基之间可能如何通过 $\subseteq$ 相联系。
        设 \( U,W \) 是某个向量空间的子空间，且 \( U\subseteq W \)。
        是否存在 \( U \) 的基 \( B_U \) 与 \( W \) 的基 \( B_W \)，使得 \( B_U\subseteq B_W \)？
        这样的基是否必定存在？

        对 \( U \) 的任意一个基 \( B_U \)，是否必有 \( W \) 的某个基 \( B_W \)，使得 \( B_U\subseteq B_W \)？

        对 \( W \) 的任意一个基 \( B_W \)，是否必有 \( U \) 的某个基 \( B_U \)，使得 \( B_U\subseteq B_W \)？

        对 \( U \) 与 \( W \) 的任意两个基 \( B_U, B_W \)，是否必有 \( B_U \) 是 \( B_W \) 的子集？
      \partsitem 基的 $\cap$ 是基吗？
        对哪个空间而言？
      \partsitem 基的 $\cup$ 是基吗？
        对哪个空间而言？
      \partsitem 那么补运算呢？
     \end{exparts}
     （\textit{提示。}用 \( \Re^3 \) 的某些子空间来检验你提出的任何猜想。）
     \begin{answer}
       首先注意，一个集合是某个空间的基当且仅当它是线性无关的，因为此时它是自身张成空间的基。
       \begin{exparts}
         \partsitem 第二段中问题的答案是“是的”
           （这意味着第一段中两个问题的答案也都是“是的”）。
           若 \( B_U \) 是 \( U \) 的一个基，则 \( B_U \) 是 \( W \) 的一个线性无关子集。
           用\nearbycorollary{cor:LIExpBas}把它扩充为 \( W \) 的一个基。
           那就是所求的 \( B_W \)。

           第三段中问题的答案是“不是”，这意味着第四段问题的答案是“不是”。
           下面举例说明：超空间的一个基，可以没有任何子基构成子空间的基——在 \( W=\Re^2 \) 中，考虑标准基 \( \stdbasis_2 \)。
           $\stdbasis_2$ 的任何子基都不能构成 $\Re^2$ 中直线 \( y=x \) 这个子空间 \( U \) 的基。
         \partsitem 它是一个基（是其张成空间的基），因为线性无关集的交是线性无关的（交是其中每一个线性无关集的子集）。

           然而，它并不是这些空间的交的基。
           例如，下面是 \( \Re^2 \) 的基：
            \begin{equation*}
              B_1=\sequence{\colvec[r]{1 \\ 0},\colvec[r]{0 \\ 1}}
              \quad\text{和}\quad
              B_2=\sequence[r]{\colvec{2 \\ 0},\colvec[r]{0 \\ 2}}
            \end{equation*}
            而 \( \Re^2\intersection\Re^2=\Re^2 \)，但 \( B_1\intersection B_2 \) 为空。
            我们只能说，基的 $\cap$ 是这些空间的交的某个子集的基。
         \partsitem 基的 $\cup$ 未必是基：在 \( \Re^2 \) 中
            \begin{equation*}
              B_1=\sequence{\colvec[r]{1 \\ 0},\colvec[r]{1 \\ 1}}
              \quad\text{和}\quad
              B_2=\sequence{\colvec[r]{1 \\ 0},\colvec[r]{0 \\ 2}}
            \end{equation*}
            \( B_1\union B_2 \) 不是线性无关的。
            两个基的 $\cup$ 为基的一个充要条件是
            \begin{equation*}
              B_1\union B_2 \text{ 线性无关 }
              \quad\iff\quad
              \spanof{B_1\intersection B_2}=\spanof{B_1}\intersection
                                             \spanof{B_2}
            \end{equation*}
            证明它足够容易（但应用起来或许很难）。
          \partsitem 一个基的补不可能是一个基，因为它包含零向量。
        \end{exparts}  
      \end{answer}
   \recommended \item 
     考虑“维数”与“子集”如何相互作用。
     设 \( U \) 与 \( W \) 是某个向量空间的子空间，且 \( U\subseteq W \)。
    \begin{exparts}
      \partsitem 证明 \( \dim (U)\leq\dim (W) \)。
      \partsitem 证明维数相等当且仅当 \( U=W \)。
      \partsitem 证明当它们是无限维的时候，上一小题不成立。
    \end{exparts}
    \begin{answer}
       \begin{exparts}
          \partsitem \( U \) 的一个基是 \( W \) 中的线性无关集，因此可以用\nearbycorollary{cor:LIExpBas}把它扩充为 \( W \) 的一个基。
            第二个基的成员至少与第一个基一样多。
          \partsitem 一个方向是清楚的：若 \( V=W \)，则二者维数相同。
            反之，设 \( B_U \) 是 \( U \) 的一个基。
            它是 \( W \) 的一个线性无关子集，因此可以扩充为 \( W \) 的一个基。
            若 \( \dim(U)=\dim(W) \)，则 \( W \) 的这个基的成员不会比 \( B_U \) 多，于是它就等于 \( B_U \)。
            由于 \( U \) 与 \( W \) 有相同的基，所以它们相等。
          \partsitem 设 \( W \) 是次数有限的多项式构成的空间，\( U \) 是只含偶数次幂项的多项式构成的子空间。
             \begin{equation*}
                U=\set{a_0+a_1x^2+a_2x^4+\dots+a_nx^{2n}\suchthat
                           a_0,\ldots,a_n\in\Re} 
             \end{equation*}
            两个空间都是无限维的，但 \( U \) 是真子空间。
        \end{exparts}  
      \end{answer}
    \item 下面给出本节主要结果
      \nearbytheorem{th:AllBasesSameSize}的另一种证明。
      先是一个例子，然后是一条引理，最后是定理。
      \begin{exparts}
        \partsitem 将这个来自 $\Re^3$ 的向量
          表示为该基的成员的线性组合。
         \begin{equation*}
           B=\sequence{\colvec[r]{1 \\ 0  \\ 0},
             \colvec[r]{1 \\ 1 \\ 0},
             \colvec[r]{0 \\ 0 \\ 2}}
           \qquad
           \vec{v}=\colvec[r]{1 \\ 2 \\ 0}
         \end{equation*}
        \partsitem 在该组合中取一个系数非零的基向量。
          通过用那个基向量替换~$\vec{v}$来改动~$B$，得到新的序列~$\hat{B}$。
          检验$\hat{B}$也是~$\Re^3$的一个基。
        \partsitem （交换引理）
          设
          \( B=\sequence{\vec{\beta}_1,\dots,\vec{\beta}_n} \) 是一个向量空间的基，且对向量 \( \vec{v} \)
          而言，关系式 \( \vec{v}=\lincombo{c}{\vec{\beta}} \)
          中有 \( c_i\neq 0 \)。
          证明用 \( \vec{v} \) 替换 \( \vec{\beta}_i \)
          后得到该空间的另一个基。
        \partsitem 利用它，结合归纳法，证明
          \nearbytheorem{th:AllBasesSameSize}。
      \end{exparts}
      \begin{answer}
      \begin{exparts}
        \partsitem
          $\colvec[r]{1 \\ 2 \\ 0}
             =(-1)\cdot\colvec[r]{1 \\ 0  \\ 0}
              +2\colvec[r]{1 \\ 1 \\ 0}
              +0\cdot\colvec[r]{0 \\ 0 \\ 2}$
         \partsitem 有两个基向量带有非零系数。
            比如我们可以取第一个。
            \begin{equation*}
              \hat{B}=\sequence{\colvec[r]{1 \\ 2  \\ 0},
                \colvec[r]{1 \\ 1 \\ 0},
                \colvec[r]{0 \\ 0 \\ 2}}
            \end{equation*}
            验证它是 \( \Re^3 \) 的另一个基是常规的。
         \partsitem  
          称交换的结果为
          \( \hat{B}=\sequence{\vec{\beta}_1,\dots,\vec{\beta}_{i-1},\vec{v},
            \vec{\beta}_{i+1},\dots,\vec{\beta}_n}  \)。
          我们先证明 $\hat{B}$ 是线性无关的。
          $\hat{B}$ 的成员之间的任一关系
          \( d_1\vec{\beta}_1+\dots+d_i\vec{v}+\dots+d_n\vec{\beta}_n=\zero \)，
          在代入 $\vec{v}$ 之后，
          \begin{equation*}
          d_1\vec{\beta}_1+\dots
          +d_i\cdot(c_1\vec{\beta}_1+\dots+c_i\vec{\beta}_i+\dots+c_n\vec{\beta}_n)
          +\dots+d_n\vec{\beta}_n
          =\zero
          \tag*{($*$)}\end{equation*}
          就给出 $B$ 的成员之间的一个线性关系。
          基 $B$ 是线性无关的，所以 $\vec{\beta}_i$ 的系数 $d_ic_i$ 为零。
          由于我们假设了 $c_i$ 非零，故 $d_i=0$。
          把它用于方程~$(*)$ 可知，其余的 $d$ 也全为零。
          因此 $\hat{B}$ 是线性无关的。

          最后我们证明 $\hat{B}$ 与 $B$ 有相同的张成。
          这个论证的一半，即
          $\spanof{\smash{\hat{B}}}\subseteq\spanof{B}$，
          是容易的；我们可以把 $\spanof{\smash{\hat{B}}}$ 的任一成员
          $d_1\vec{\beta}_1+\dots+d_i\vec{v}+\dots+d_n\vec{\beta}_n$
          写成
          $
          d_1\vec{\beta}_1+\dots
          +d_i\cdot(c_1\vec{\beta}_1+\dots+c_n\vec{\beta}_n)
          +\dots+d_n\vec{\beta}_n
          $，
          它是 $B$ 的成员的线性组合的线性组合，从而属于 $\spanof{B}$。
          对于论证的另一半 $\spanof{B}\subseteq\spanof{\smash{\hat{B}}}$，回忆一下：若
          $\vec{v}=c_1\vec{\beta}_1+\dots+c_n\vec{\beta}_n$ 且 $c_i\neq 0$，
          则可以把该等式改写为
          $\vec{\beta}_i=(-c_1/c_i)\vec{\beta}_1+\dots+(1/c_i)\vec{v}+\dots
          +(-c_n/c_i)\vec{\beta}_n$。
          现在，考虑 $\spanof{B}$ 的任一成员
          $d_1\vec{\beta}_1+\dots+d_i\vec{\beta}_i+\dots+d_n\vec{\beta}_n$，
          把 $\vec{\beta}_i$ 替换为它关于 $\hat{B}$ 的成员的线性组合表达式，并像本论证的前一半那样看出，
          结果是 $\hat{B}$ 的成员的线性组合的线性组合，
          从而属于 $\spanof{\smash{\hat{B}}}$。
         \partsitem 取定一个至少有一个有限基的向量空间。
           从这个空间的所有基中，选出一个规模最小的
           \( B=\sequence{\vec{\beta}_1,\dots,\vec{\beta}_n} \)。
           我们将证明任何其他的基
           \( D=\smash{\sequence{\vec{\delta}_1,\vec{\delta}_2,\ldots}} \)
           也有相同的成员数 $n$。
           因为 \( B \) 规模最小，\( D \) 有不少于 \( n \)~个向量。
           我们将论证它不可能有多于 \( n \) 个向量。

           基 \( B \) 张成该空间，而 \( \vec{\delta}_1 \)
           在该空间中，
           所以 \( \vec{\delta}_1 \) 是 \( B \) 的元素的一个非平凡线性组合。
           由交换引理，可以把 \( B \) 中的一个向量换成
           \( \vec{\delta}_1 \)，得到基 \( B_1 \)，
           其中一个元素是
           \( \vec{\delta}_1 \)，
           而其余 \( n-1 \) 个元素都是 \( \vec{\beta} \)。

           上一段构成归纳论证的基础步骤。
           归纳步骤从一个基 \( B_k \)（其中 \( 1\leq k<n \)）开始，
           它包含 \( D \) 的 \( k \) 个成员与 \( B \) 的 \( n-k \) 个成员。
           我们知道 \( D \) 至少有 \( n \) 个成员，所以存在
           \( \vec{\delta}_{k+1} \)。
           把它表示为 \( B_k \) 的元素的线性组合。
           关键点：在那个表示中，
           至少有一个非零标量必须与某个 \( \vec{\beta}_i \) 相关联，
           否则该表示就会是线性无关集 \( D \) 的元素之间的一个
           非平凡线性关系。
           用 \( \vec{\delta}_{k+1} \) 换掉 \( \vec{\beta}_i \)，
           得到新的基
           \( B_{k+1} \)，与前面的基 \( B_k \) 相比，它多一个 \( \vec{\delta} \)、少一个
           \( \vec{\beta} \)。

           重复这一步，直到没有 \( \vec{\beta} \) 剩下，
           于是 \( B_n \) 包含  
           $\vec{\delta}_1,\dots,\vec{\delta}_n$。
           现在，\( D \) 不可能有多于这 \( n \) 个的向量，因为
           剩下的任何 \( \vec{\delta}_{n+1} \) 都会落在
           \( B_n \)（因为它是一个基）的张成中，
           从而是其余 $\vec{\delta}$ 的线性组合，
           这与 $D$ 线性无关矛盾。
      \end{exparts}
      \end{answer}
  \puzzle \item 
    \cite{Sheffer}
    一个图书馆有 $n$ 本书和 $n+1$ 位订户。
    每位订户都从该图书馆至少读过一本书。
    证明必存在两个不相交的订户集合，他们读过的书完全相同
    （也就是说，每个集合中的订户读过的书的并集是相同的）。
    \begin{answer}
      （这个解答来自网站上 Yuzhou Gu 的一条评论。）
      给每个人指定一个向量 $\vec{v}_i\in\R^n$，
      若此人读过该书，则相应分量为 $1$，否则为 $0$。
      由于这些向量线性相关，我们将得到一个形如 $\sum a_i v_i = 0$ 的等式。
      那么，使 $a_i>0$ 的那些 $i$ 构成的集合与使
      $a_i<0$ 的那些 $i$ 构成的集合，就是两个不相交的集合，且读过的书的并集相同。
    \end{answer}
  \puzzle \item 
    \cite{Wohascum47}
    对 $\Re^n$ 中任意向量 $\vec{v}$ 以及数 $1$、$2$、\ldots、$n$ 的任意置换 $\sigma$
    （也就是说，$\sigma$ 是把这些数重排成新次序的一个排列），
    定义 $\sigma(\vec{v})$ 为分量为
    $v_{\sigma(1)}$、$v_{\sigma(2)}$、\ldots{}、$v_{\sigma(n)}$ 的向量
    （其中 $\sigma(1)$ 是重排后的第一个数，等等）。
    现在取定 $\vec{v}$，令 $V$ 为
    $\set{\sigma(\vec{v})\suchthat \mbox{$\sigma$ 置换 $1$, \ldots, $n$}}$ 的张成。
    $V$ 的维数有哪些可能？
    \begin{answer}
      $V$ 的维数的可能是 $0$、$1$、$n-1$
      和 $n$。

      为看到这一点，先考虑 $\vec{v}$ 的所有坐标都相等的情形。
      \begin{equation*}
         \vec{v}=\colvec{z \\ z \\ \vdots \\ z}
      \end{equation*}
      那么对每个置换 $\sigma$ 都有 $\sigma(\vec{v})=\vec{v}$，
      所以 $V$ 就是 $\vec{v}$ 的张成，
      依 $\vec{v}$ 是否为 $\zero$ 而定，其维数为 $0$ 或 $1$。

      现在假设 $\vec{v}$ 的坐标不全相等；设 $x$
      与 $y$（$x\neq y$）是 $\vec{v}$ 的坐标中的两个。
      那么我们可以找到置换 $\sigma_1$ 与 $\sigma_2$，使得
      \begin{equation*}
        \sigma_1(\vec{v})=\colvec{x \\ y \\ a_3 \\ \vdots \\ a_n}
        \quad\text{和}\quad
        \sigma_2(\vec{v})=\colvec{y \\ x \\ a_3 \\ \vdots \\ a_n}
      \end{equation*}
      其中 $a_3,\ldots,a_n\in\Re$。
      因此，
      \begin{equation*}
        \frac{1}{y-x}\bigl(\sigma_1(\vec{v})-\sigma_2(\vec{v})\bigr)=
        \colvec[r]{-1 \\ 1 \\ 0 \\ \vdotswithin{-1} \\ 0}
      \end{equation*}
      属于 $V$。
      也就是说，$\vec{e}_2-\vec{e}_1\in V$，其中 $\vec{e}_1$、$\vec{e}_2$、
      \ldots、$\vec{e}_n$ 是 $\Re^n$ 的标准基。
      类似地，$\vec{e}_3-\vec{e}_2$、\ldots、$\vec{e}_n-\vec{e}_1$
      都在 $V$ 中。
      容易看出，向量
      $\vec{e}_2-\vec{e}_1$、$\vec{e}_3-\vec{e}_2$、\ldots、
      $\vec{e}_n-\vec{e}_1$ 是线性无关的（也就是说，构成一个线性无关集），
      所以 $\dim V\geq n-1$。

      最后，我们可以写出
      \begin{align*}
        \vec{v} &=x_1\vec{e}_1+x_2\vec{e}_2+\dots+x_n\vec{e}_n  \\
                &=(x_1+x_2+\dots+x_n)\vec{e}_1+x_2(\vec{e}_2-\vec{e}_1)+\dots
                   +x_n(\vec{e}_n-\vec{e}_1)
      \end{align*}
      这表明，若 $x_1+x_2+\dots+x_n=0$，则 $\vec{v}$ 落在
      $\vec{e}_2-\vec{e}_1$、\ldots、$\vec{e_n}-\vec{e}_1$ 的张成中
      （也就是说，落在这些向量构成的集合的张成中）；
      类似地，每个 $\sigma(\vec{v})$
      也在这个张成中，所以 $V$ 就等于这个张成，且 $\dim V=n-1$。
      另一方面，若 $x_1+x_2+\cdots+x_n\neq 0$，则上面的等式
      表明 $\vec{e}_1\in V$，从而 $\vec{e}_1,\dots,\vec{e}_n\in V$，
      所以 $V=\Re^n$ 且 $\dim V=n$。
    \end{answer}
\end{exercises}
\index{basis|)}



























